\scanpage{13}
% source lines 1-3
\[
A(K_{\mathfrak p})\simeq K_{\mathfrak p}^*
\]
de sorte qu'on a des épimorphismes
\[
{\unsure{\varphi|\phi}}_{\mathfrak p}:K_{\mathfrak p}^*\longrightarrow G_{\mathfrak p}^{d}
\]

% source lines 4-5, circled 5 degree
\textcircled{$5^\circ$} \uline{Ces applications sont précisément les hom. de réciprocité, \unsure{donnés|données} par les symboles de \deleted{\unsure{Art|Artin}} \inserted{\unsure{restes}} \unsure{normiques}.}

% source lines 6-12
Comme les deux applications coïncident manifestement aux places non ramifiées, l'assertion (5) équivaut au fait que ces homs $\varphi$ \inserted{\unsure{continus|continu}} proviennent bien d'un hom unique $C(K)\longrightarrow G_{\mathfrak p}^{d}$, i.e. \unsure{ils|il} \uline{signifient} \uline{essentiellement} que l'on a la formule
\[
\boxed{\prod_{\mathfrak p}\varphi_{\mathfrak p}(x_{\mathfrak p})=1\qquad\text{si }x\in K^*}
\]

% source lines 14-16
(Formule qui \deleted{\unsure{est loin d'être}} équivaut à la partie difficile de la théorie \unsure{actuelle} des corps de classes pour les corps de \unsure{nbs|nbres}).
% archive sheet number: 12
